% !TeX root = ../tesis.tex

% -----------------------------------------------------------------------
\section{Datos Abiertos y \gis{} Participativo como Infraestructura Metodológica}
\label{sec:sig-datos-abiertos}
% -----------------------------------------------------------------------

La viabilidad de una metodología replicable para el diagnóstico topológico de
redes de transporte metropolitano descansa, en gran medida, en la disponibilidad
de datos geoespaciales abiertos y estandarizados. La presente investigación
utiliza dos tipos de fuentes complementarias: los feeds \gtfs{} oficiales de los
sistemas de transporte masivo de la \zmvm{} y
\textbf{la cartografía de base derivada de iniciativas de información geográfica voluntaria (\textsc{vgi}}
, por sus siglas en inglés). Esta sección establece el sustento teórico que
justifica la pertinencia y suficiencia de ambas fuentes para el análisis a
escala metropolitana, así como el argumento de equidad que fundamenta el
desarrollo de herramientas de código abierto accesibles para planificadores
urbanos.

% -----------------------------------------------------------------------
\subsection{VGI y cobertura de datos a escala metropolitana}
\label{subsec:vgi-cobertura}
% -----------------------------------------------------------------------

La validez del análisis topológico propuesto descansa, en parte, en la calidad y
completitud de los datos geoespaciales de base. \citet{haklay2010good} demostró,
mediante una comparación sistemática entre datos de OpenStreetMap y los del
\textit{Ordnance Survey} del Reino Unido, que la información \textsc{vgi}
alcanza niveles de precisión posicional superiores al 80\% en áreas urbanas
densas, resultado atribuible a la alta concentración de contribuidores activos
en dichos entornos. Este umbral de calidad es comparable al de fuentes
cartográficas institucionales, lo que posiciona a los datos colaborativos como
una alternativa técnicamente válida para el modelado de redes en ciudades con
cobertura de contribuidores suficiente.

Complementariamente, \citet{barrington2017world} evidenciaron que la red vial de
OpenStreetMap supera el 80\% de completitud a escala global, con niveles aún más
elevados en las zonas metropolitanas de países de ingresos medios y altos. Su
análisis también demostró que las ciudades latinoamericanas de escala comparable
a la \zmvm{} presentan coberturas superiores al promedio global, resultado del
crecimiento sostenido de comunidades locales de cartografía colaborativa. Estos
hallazgos sustentan la pertinencia de emplear datos \gis{} de fuentes abiertas
---en combinación con feeds \gtfs{} oficiales--- como insumo primario para el
modelado de la red multimodal de la \zmvm{}.

% -----------------------------------------------------------------------
\subsection{Datos abiertos, comunidad y equidad de acceso}
\label{subsec:datos-comunidad-equidad}
% -----------------------------------------------------------------------

La disponibilidad de datos no garantiza, por sí sola, su aprovechamiento
equitativo. \citet{gurstein2011open} argumenta que la diferencia entre datos
abiertos y datos verdaderamente útiles reside en la existencia de herramientas
que traduzcan la información bruta en conocimiento accionable para actores con
capacidades técnicas heterogéneas; sin dicha mediación, los datos abiertos
empoderan únicamente a quienes ya cuentan con los recursos computacionales y el
capital humano para procesarlos, reproduciendo en el plano digital las mismas
desigualdades de acceso que caracterizan la movilidad metropolitana.

Por su parte, \citet{neis2012street} documentaron cómo la colaboración
distribuida en OpenStreetMap generó, entre 2007 y 2011, una red vial completa y
detallada para todo el territorio alemán, demostrando que los datos \gis{} de
base colaborativa pueden alcanzar la escala y densidad requeridas para el
análisis urbano de precisión en plazos relativamente cortos. La evolución
acelerada de la cobertura cartográfica comunitaria valida que las ciudades que
actualmente disponen de datos incompletos pueden incorporarse progresivamente a
este tipo de análisis conforme la comunidad local de contribuidores se
consolida. En conjunto, estas perspectivas justifican la necesidad de
desarrollar herramientas de código abierto que medien entre los datos
disponibles y los planificadores urbanos que carecen de formación especializada
en programación o análisis de redes complejas.

% -----------------------------------------------------------------------
\subsection{De los datos abiertos a la herramienta replicable}
\label{subsec:datos-herramienta-replicable}
% -----------------------------------------------------------------------

La cadena metodológica propuesta en esta investigación ---\textit{Apimetro} para
la adquisición y normalización de feeds \gtfs{}, \textit{VFTModel} para el
cálculo de indicadores topológicos y \textit{Transport-gis} para la
visualización cartográfica de resultados--- representa una respuesta concreta a
la brecha identificada por \citet{gurstein2011open}: tres motores de código
abierto que transforman datos estandarizados en diagnósticos cuantitativos
comprensibles para urbanistas sin especialización en ciencias de datos. Esta
arquitectura metodológica es, por diseño, independiente de licencias comerciales
y reproducible en cualquier área metropolitana que disponga de feeds \gtfs{} y
cartografía \gis{} de base, condición que, conforme a los hallazgos de
\citet{barrington2017world}, se cumple en la gran mayoría de las ciudades
latinoamericanas de escala comparable a la \zmvm{}.

El argumento de replicabilidad no es secundario, sino central al objetivo
general de la investigación: demostrar que un planificador urbano con acceso a
datos abiertos puede, mediante esta metodología, obtener un diagnóstico
topológico riguroso de su red de transporte y evaluar escenarios alternativos de
intervención sin depender de software propietario ni de equipos técnicos
especializados. Los referentes de calidad de datos
\citep{haklay2010good, barrington2017world} y los principios de equidad en el
acceso al conocimiento \citep{gurstein2011open, neis2012street} constituyen, en
este sentido, el fundamento teórico que articula la dimensión metodológica de la
tesis con su dimensión urbanística.

